\documentclass[a4paper]{article}
\usepackage{amsmath}
\usepackage{amssymb}
\usepackage{geometry}
\usepackage{hyperref}

\geometry{margin=25mm}

\title{The Electromagnetic Host of the Slow Charged Entrance:\\
A Restoring Seat in an Octahedral Cage}

\author{https://github.com/hypernumbernet}
\date{\today}

\begin{document}
\maketitle

\begin{abstract}
The slow charged entrance of the nuclear channel is a path on which the relative rapidity stays below the falling threshold while the separation approaches the nuclear node, and on which the electromagnetic factor stays positive at that node. A free infall meets the threshold before the node. The nuclear force fixes the sign and does not name what holds the rapidity down. Along one axis the holding is electromagnetic. It is a cage of like ions.

A regular octahedron of positive ions, each of charge at least one, holds a further positive ion at its center and pushes an electron off that center. The electrons of the metal are free. Screening is not introduced. Two protons on an axis of the same octahedron are a different seat. With the cage radius equal to twice the electromagnetic length, the component that increases their separation is strictly negative at separation twice that length, for every such ion charge. Just outside the first charge node the same component tends to infinity: the repulsion of the partner diverges, and every ion stays off the node. The component therefore vanishes at a separation strictly between the node and that diameter. Just outside the node the pair is driven outward, and at the diameter it is driven inward. The cage does not carry the pair through the node.

The seat withholds the separation from a free electromagnetic infall, and it is a host of that keeping. A deuteron carries the same unit charge, so the component is the component of a deuteron pair. The seat does not set the ratio of the electromagnetic length to the nuclear length. If the charge node lies outside the nuclear node, the seat stands outside an infinite electromagnetic wall, and the classical path on the axis does not reach the nuclear skin. The rate, an excess of heat, and the choice among experiments remain underived. The coordination saddles of a metallic bond are not this seat.
\end{abstract}

\section{The cage}
\label{sec:cage}

The nuclear companion fixes the slow charged entrance and leaves its host unnamed~\cite{dst-strong}. Below the rapidity threshold \(\varphi_*\) the nuclear term opposes the Coulomb term, and next to the threshold it exceeds every finite ratio of the two prefactors, so long as the electromagnetic factor stays positive. Above the threshold the nuclear term adds to the Coulomb repulsion. A free infall still meets the threshold before the node, so the entrance is a path on which the rapidity is held down. When the kinetic energy is at most \(10^{-10}\) of the rest energy, the fractional widening of the radius is less than \(5\times 10^{-6}\), and a nucleon or a deuteron at room temperature is of this kind. The present note names the host along one axis. The naming uses the electromagnetic channel already fixed for charges~\cite{dst-shell}. No second coupling is introduced.

On the equal-scale locus the force that increases the separation of two charges is
\begin{equation}
\label{eq:force}
F(r)=\frac{k q_1 q_2}{\gamma_s(r)\,r^{2}},
\qquad
k=\frac{1}{4\pi\epsilon_0},
\end{equation}
with \(\gamma_s(x)=\cosh x\cos x-\sinh x\sin x\) and \(x=\ell/(2r)\). Like charges repel where \(\gamma_s>0\) and attract where \(\gamma_s<0\). The first node \(x_1\) is the unique root of \(\tan x=\coth x\) in \((\pi/4,\pi/2)\), at \(r_1/\ell=0.53330\). The outer well is \(r>r_1\), where \(\gamma_s>0\).

A regular octahedron, a cube, and a cuboctahedron of positive ions hold a further positive ion at the center and push an electron off that center. One or two of the ions may be absent. The force at the center vanishes by the symmetry, and every displacement of the positive ion is restored, while every displacement of an electron is pushed further. Under the Coulomb law both rates are zero: that zero is Earnshaw's theorem, and the signed factor removes it in opposite directions for the two signs of the charge. The same holding is the holding of every cubic truncation, out to any finite number of shells, while every charge is positive and every phase lies below \(\pi\) and off the first node. A metal in the outer well is a held arrangement of ions. The electrons have no seat among them. Screening is not introduced, and each ion remains a point charge. No infinite lattice sum is taken.

The nearest positive ions about an octahedral hole of the face-centered cubic arrangement form such an octahedron, and the next shell is a cube. Palladium and nickel are of that arrangement. The name records the geometry of the hole. Fleischmann and Pons placed deuterium in palladium and asked whether a nuclear entrance occurred there~\cite{fleischmann1989}. The seat below does not decide what that arrangement produced.

\section{The axial component}
\label{sec:axis}

The centered ion is one charge, restored in every direction. Two protons on an axis of the octahedron are a pair, and the component computed here is the component along that axis. Place ions of charge \(Ze\) at \((\pm a,0,0)\), \((0,\pm a,0)\), and \((0,0,\pm a)\), and protons of charge \(+e\) at \((\pm t,0,0)\), with \(0<t<a\). Every ion--proton separation and the proton--proton separation are taken off every node. The component that increases the separation is
\begin{equation}
\label{eq:axial}
F_{\parallel}(t)
=
-Z F(a-t)+Z F(a+t)
+4Z\frac{t}{R}F(R)
+F(2t),
\end{equation}
where \(R=\sqrt{a^{2}+t^{2}}\) and \(F\) is~\eqref{eq:force} for two unit charges. The first term is the nearer axial ion. On the outer well it is a repulsion, and the direction from that ion toward the proton points inward, so the term decreases the separation. The second term is the farther axial ion, and it increases the separation. The third is the four equatorial ions: each repulsion has an outward axial projection \(t/R\). The last term is the partner. On the outer well, \(F_{\parallel}\) is a contest among repulsions.

The coordination saddles of a metallic bond sit at a reduced electron content, \(\nu=2/z\), for \(z=8\) and \(z=12\). Each is a saddle of that content. Neither is the interstitial octahedron, and neither is~\eqref{eq:axial}.

\section{Outward at the node, inward at the diameter}
\label{sec:seat}

Take the cage radius to be twice the electromagnetic length, \(a=2\ell\), and the half-separation to be \(t=\ell\). The separation of the protons is then the diameter \(2\ell\). The nearer ion stands at \(\ell\), the farther at \(3\ell\), each equatorial ion at \(\sqrt{5}\,\ell\), and the partner at \(2\ell\). Each of these distances lies strictly outside the first node, in the outer well. For every ion charge \(Z\ge 1\),
\begin{equation}
\label{eq:inward}
F_{\parallel}(\ell)<0.
\end{equation}
The nearer ion outweighs the farther ion, the four equatorial ions, and the partner together. Its phase is the largest of the four phases, so \(1/\gamma_s\) amplifies it more than it amplifies the others. The cage contribution at this separation is itself inward, and a larger \(Z\) adds to that inward side. The partner contribution does not depend on \(Z\).

The partner term is different in kind once the separation falls. As the separation decreases to the first charge node, \(2t\downarrow r_1\), the phase of the partner increases to \(x_1\) and \(F(2t)\) tends to \(+\infty\). At cage radius \(2\ell\) the six ions remain at distances bounded away from the node: the nearest ion approaches \(2\ell-r_1/2\), which is still outside \(r_1\), and the other five are farther. Their contribution tends to a finite value. Therefore
\begin{equation}
\label{eq:diverge}
F_{\parallel}(t)\to +\infty
\qquad\text{as}\qquad
2t\downarrow r_1.
\end{equation}
The axial component is outward in a punctured right neighborhood of the node and inward at the diameter \(2\ell\). It is continuous on the intervening separations, all of which stay in the outer well. It vanishes at some separation strictly between the node and the diameter:
\begin{equation}
\label{eq:balance}
r_1<s<2\ell,
\qquad
F_{\parallel}(s/2)=0.
\end{equation}
From the node end the partner drives the separation outward, back into the open interval. From the diameter the cage drives it inward, back into the same interval. The cage does not carry the pair through the node.

The divergence that stops the approach sits at the finite radius \(r_1\), where \(\gamma_s\) vanishes. Under the Coulomb law the same partner term diverges only at coincidence, and the three-direction sum on a single charge is zero, so the cage does not hold one proton at its center. The signed factor does both. It holds one positive ion at the center, and it stops a pair at a finite node. The zero~\eqref{eq:balance} is a vanishing of the axial component between those two signs. The note does not count how many such zeros the interval contains.

A deuteron carries charge \(+e\). The component~\eqref{eq:axial} is therefore the axial component of a deuteron pair in the same cage. The nuclear comparison for that pair remains the charged entrance of the companion: the electromagnetic factor positive at the nuclear node, and the rapidity below \(\varphi_*\).

\section{What the seat leaves open}
\label{sec:open}

The seat withholds the separation from a free electromagnetic infall. On this axis the separation is driven into the interval \((r_1,2\ell)\), and it is not left to fall through the charge node. The relative rapidity of a free infall is the rapidity of a separation that is falling toward the node. The seat keeps the separation off that fall, and in that sense it is a host of the keeping that the nuclear companion asked for. Room temperature remains the class already fixed there.

The two lengths stay independent. The seat is a statement in the electromagnetic length \(\ell\). The nuclear entrance is a statement in the nuclear length: the electromagnetic factor positive at the nuclear node, which is the ordering in which that node is the first wall met from large separation, and the rapidity below \(\varphi_*\). The cage does not set the ratio of the lengths. If instead the charge node lies outside the nuclear node, the balance~\eqref{eq:balance} stands outside an infinite electromagnetic wall. The classical path on the axis does not reach the nuclear skin. In either ordering the seat does not derive a rate, an excess of heat, or a choice among the experiments that have asked whether a cold charged entrance occurs. The sign along the axis is what it fixes.

\begin{thebibliography}{9}

\bibitem{dst-strong}
Anonymous,
``The Strong Force as a Repulsive Layer of Proximity Gravity: Proton Chaos, Magic Numbers, Slow Neutrons, and Cold Fusion,''
companion note.

\bibitem{dst-shell}
Anonymous,
``Electron Orbits from the Signed Electrostatic Force: Covalent Bonds, Metallic Bonds, and Crystals,''
companion note.

\bibitem{fleischmann1989}
M.~Fleischmann and S.~Pons,
``Electrochemically induced nuclear fusion of deuterium,''
\textit{J.~Electroanal.~Chem.} \textbf{261}, 301 (1989).

\end{thebibliography}

\end{document}
